\chapter{Discussion and Future Work}
\label{chap:conclusion}
{%

\section{Theoretical Implications}
\label{sec:theoretical-implications}

This thesis formalises the Cloud GPU Rendering Scheduling Problem (CGRSP) as an online bi-objective scheduling problem on heterogeneous, stochastically available GPU fleets. The gap it fills is specific: cloud GPU rental markets (RunPod, Lambda Labs, CoreWeave) expose fine-grained billing (per second on RunPod), qualitative node availability, and mixed deadline constraints that reserved-cluster schedulers do not handle, yet no prior scheduling work addresses all three simultaneously.

\textbf{Problem formulation.} Chapter~\ref{chap:formulation} develops this into notation, decision variables, and constraints for scheduling under stochastic provisioning delays. The three-state availability model (unavailable / idle / running; the unavailable state does not occur on the on-demand platform studied) captures the consumer's limited observability of cloud node state, without knowledge of the underlying provisioning dynamics, and is sufficient for practical policy construction, while the bi-objective function explicitly balances QoS (waiting time, deadline miss rate) against operational cost.

\textbf{Simulator construction.} Chapter~\ref{chap:system} describes the discrete-event simulator, calibrated against 6,234 real RunPod rendering jobs. The simulator reproduces the empirical execution-time distribution and adds two modelled components: Poisson job arrivals at day-type rates (the hectic and surge loads being stress assumptions above the busiest recorded day), and stochastic, stock-status-driven provisioning delays whose stock-status probabilities are assumed rather than measured. Execution times are sampled from per-stratum log-normal service distributions whose means equal the per-stratum means of those jobs for the nine well-sampled strata (the three sparse strata use relative-performance estimates, see the limitations below) and whose log-space dispersions are fitted from the same data. Simulation is used because live A/B testing on a production cluster is not feasible at this evaluation scale.

\textbf{Algorithm comparison.} Four experimental phases with heterogeneous deadlines (20\% tight 1h / 80\% loose 8h) support the following conclusions.

SPT achieves a lowest-tier wait under capacity saturation ($\rho>1$), with no significant paired wait difference from RH, though EDF now achieves significantly lower miss. Across $n=30$ seeds on the assumed hectic-day stress workload ($\sim$950 jobs/day, against a busiest recorded day of 415; $C_{\max}=5$), it records a lowest-tier wait (\SPTwait~min, $d=\SPTFIFOwaitD$ vs FIFO, $p<0.001$; with no significant paired wait difference from RH's \RHWait~min, $p=\RHSPTwaitP$) but a significantly higher miss rate than EDF (\SPTmiss\% vs EDF's \EDFmiss\%, $d=\SPTEDFmissD$, $p<0.001$). This finding is consistent with the classical single-machine SPT result~\cite{smith1956various} and its M/G/1 queueing counterpart~\cite{pinedo2016scheduling} in the heterogeneous multi-machine setting of this simulation: under heavy overload, shortest-job-first ordering achieves the lowest mean wait among the baseline policies, but deadline ordering alone now achieves significantly better deadline performance.
SPT+Rescue limits tight-deadline starvation at a notably lower miss rate than SPT (\RescueMiss\% vs \SPTmiss\%). The laxity-rescue mechanism achieves \RescueMiss\% miss ($d=\RescueFIFOmissD$, $p<0.001$ vs FIFO) with \RescueWait~min wait ($d=\RescueFIFOwaitD$, $p<0.001$), a conservative alternative to SPT when tight-deadline compliance is a strict operational requirement.

RH has the lowest deadline miss point estimate (\RHMiss\%, with no significant paired difference from CADR's in this sample) with lowest-tier wait. Under hectic-day saturation (Phase~3, $n=30$ seeds), RH's miss advantage over SPT+Rescue survives Holm correction ($d=\RHRescueMissD$, $p<0.001$), while its wait advantage, though numerically favourable ($d=\RHRescueWaitD$; \RHWait~min vs \RescueWait~min), has a raw $p=\RHRescueWaitP$ that does not survive the family-wise correction: RH is significantly better than SPT+Rescue on miss and comparable to it on wait, at an operational cost within roughly 1\% of SPT+Rescue. RH and CADR achieve near-identical lowest monthly miss in Phase~4 (RH \PhFourRHMiss\%, CADR \PhFourCADRMiss\%; no paired significance test is computed at this phase), though at monthly scale RH's wait no longer leads: it sits within about two minutes of SPT's, CADR's and SPT+Rescue's (\PhFourRHWait~min vs \PhFourSPTwait~min, \PhFourCADRwait~min and \PhFourRescueWait~min, no paired policy test). The forward-simulation plan keeps shortest-processing-time ordering for jobs with slack, promotes only jobs the timeline shows would otherwise miss, and demotes already-doomed jobs so they do not displace savable ones. RH's late promotion of tight jobs, which then fall into its demoted tier, results in high mean tardiness (\RHtard~min Phase~3, \PhFourRHtard~min Phase~4), now exceeding even the deadline-oblivious FIFO baseline: missed deadlines under RH are missed by substantially more than under CADR or EDF, and by more than under FIFO's blind ordering. RH's anticipation is load-gated: below saturation it reserves a slot for imminent tight (rush) arrivals and dispatches waiting tight jobs eagerly, cutting surge-day miss from \ResSurgeOffMiss\% to \ResSurgeOnMiss\% (tight-deadline miss from \ResSurgeOffTight\% to \ResSurgeOnTight\%); the offered-load gate disables reservation under saturation, so the hectic-day headline (\ResHecticGateMiss\% miss) is unchanged (Section~\ref{ssec:anticipation-reservation-sensitivity}). The monthly headline is not: a month mixes surge, normal and quiet days on which the gate leaves anticipation live, and the paired Phase~4 ablation over the same ten replications shows that switching anticipation off raises monthly miss above \PhFourRHMiss\% and monthly mean tardiness above \PhFourRHtard~min ($p<0.001$ on both), while trimming a little from monthly wait. RH's monthly result therefore depends on anticipation, and its hectic-day result does not.

CADR achieves near-RH miss while substantially reducing tardiness. In Phase~3 it records \CADRmiss\% miss (vs RH \RHMiss\%) at \CADRtard~min mean tardiness (vs RH \RHtard~min), a reduction of \CADRtardReduction\%; in Phase~4 monthly operation, \PhFourCADRMiss\% miss at \PhFourCADRTard~min tardiness (vs RH \PhFourRHtard~min). At Phase~3 hectic saturation CADR has a lower miss point estimate than SPT+Rescue (\CADRmiss\% vs \RescueMiss\%, a moderate effect with no pre-specified paired test) at essentially the same wait (\CADRwait~min vs \RescueWait~min) and comparable tardiness: the advantage over SPT+Rescue is concentrated on miss rate. The cost-aware placement component is a small refinement: the CADR-order-only ablation (\CADROOmiss\% miss, \CADROOtard~min tardiness) shows that the critical-ratio ordering drives most of CADR's advantage. Operational cost headroom is structurally small across all policies (a marginal, near-noise gap), and CADR is not the cheapest scheduler overall. The critical-ratio miss is flat for $\mathrm{CR}^*\geq 3$ (best \CRStarBestMiss\% at $\mathrm{CR}^*=3$) and rises at $\mathrm{CR}^*\leq 2$ (CADR's lead over SPT+Rescue narrows at 2 and is gone at 1.5; spread \CRStarMissSpread\,pp across the swept range).

EDF now achieves significantly lower miss than SPT, at a wait that the paired test does not separate from FIFO's. EDF's deadline-ordering achieves \EDFmiss\% miss ($d=\EDFFIFOmissD$, $p<0.001$ vs FIFO) and leaves waiting time at FIFO's level ($d=\EDFFIFOwaitD$, $p=\EDFFIFOwaitP$, not significant) rather than reducing it. Under hectic saturation, EDF is dominated by SPT on wait, but significantly beats SPT on miss ($d=\SPTEDFmissD$, $p<0.001$). EDF's mean tardiness (\EDFtard~min Phase~3, \PhFourEDFtard~min Phase~4) remains low relative to the deadline-oblivious policies, because its deadline ordering ensures even missed jobs are missed by a small margin, though at Phase~3 it no longer leads the deadline-aware group, the Adaptive hybrid posting lower tardiness there (\AdaptiveTard~min vs EDF's \EDFtard~min); at Phase~4 the Adaptive hybrid also has lower tardiness (\PhFourAdaptiveTard~min versus EDF's \PhFourEDFtard~min).

FIFO, LCF, and Balanced incur \FIFOmiss\% to \Balancedmiss\% miss. Without deadline information in the dispatch logic, tight-deadline jobs waiting behind loose-deadline jobs routinely miss their 1-hour window. The result confirms that deadline-oblivious policies are unsuitable for mixed-deadline rendering workloads.

The monthly simulation, whose hectic days carry the same assumed stress load and make up most of its volume, reproduces the saturation finding (SPT: \PhFourSPTwait~min, \PhFourSPTmiss\% miss; CADR: \PhFourCADRwait~min, \PhFourCADRMiss\% miss; SPT+Rescue: \PhFourRescueWait~min, \PhFourRescueMiss\% miss; EDF: \PhFourEDFwait~min, \PhFourEDFmiss\%; FIFO/\allowbreak LCF/\allowbreak Balanced: ${\approx}$\PhFourFIFOmiss{} to \PhFourLCFmiss\% miss). At the assumed hectic-day load, a stress load above anything in the job record (Section~\ref{ssec:platform-and-data}), the concurrency budget $C_{\max}=5$ is the binding throughput constraint, and no scheduling policy alone can overcome it; only raising the budget can. At the quiet, normal and surge loads the execution-only offered load stays below saturation, although on surge days the provisioning delays, during which a dispatched job already holds its slot, still build mean waits of \SensSrgDayWaitLo{} to \SensSrgDayWaitHi{}~min. The tardiness-objective study (Section~\ref{ssec:objective-weight-sensitivity}) shows that the per-metric leaders are stable, RH lowest on miss point estimate and CADR lowest on tardiness among the two proposed schedulers, but that the single scalarised winner depends on whether deadline violations are penalised as a binary miss indicator or as continuous tardiness; because the baseline binary miss comparison is inconclusive while tardiness separates the two sharply, CADR is recommended as the default scheduler, and RH is kept as the specialist for the scarce baseline availability regime.

\textbf{Practical guidance.} Based on these findings, the following deployment rules are recommended for cloud rendering operators:

\begin{enumerate}
  \item \textbf{Use CADR as the default scheduler.} CADR's hectic-saturation miss (\CADRmiss\%) does not differ significantly from RH's (\RHMiss\%, $p=\RHCADRmissP$), but equivalence was not established, and CADR has close monthly miss point estimates with lower tardiness, and leads RH under eased availability and the data-ordered per-type ranking: at monthly scale its miss point estimate is slightly higher (\PhFourCADRMiss\% vs RH's \PhFourRHMiss\%, no paired test at that phase), at the eased $\times 1.25$ availability scale it reaches zero misses while RH does not (\BetaUpRHmiss\%), and under the data-ordered per-type ranking, the one the pricing-record offer shares support, it leads and RH is significantly worse (\RevCADRmiss\% vs \RevRHmiss\%, $\RevRHvCADRp$) (Section~\ref{ssec:availability-model-sensitivity}). Its tardiness is substantially lower at every phase (\CADRtard~min vs RH's \RHtard~min at Phase~3, \PhFourCADRTard~min vs \PhFourRHtard~min at Phase~4), because its earlier critical-ratio promotion leaves few tight jobs to be demoted, whereas RH's misses are mostly tight jobs missed by a large margin. CADR also has a lower miss point estimate than SPT+Rescue at hectic saturation (a moderate effect, no pre-specified paired test) while essentially matching its wait. The critical-ratio miss is flat for $\mathrm{CR}^*\geq 3$ (best \CRStarBestMiss\% at $\mathrm{CR}^*=3$) and rises at $\mathrm{CR}^*\leq 2$ (CADR's lead over SPT+Rescue narrows at 2 and is gone at 1.5; spread \CRStarMissSpread\,pp). CADR is not the cheapest scheduler, but operational cost differences across policies are a marginal, near-noise gap.
  \item \textbf{Use RH as the specialist for the scarce baseline availability regime, when only the count of missed deadlines matters.} This holds at the binding default budget $C_{\max}=5$; in the simulated $C_{\max}=10$ what-if every scheduler reaches zero miss and RH has the highest wait of the seven schedulers (\CmaxTenRHwait~min against \CmaxTenOtherWaitLo{} to \CmaxTenOtherWaitHi~min for the others), so there CADR or SPT is the better choice. RH has the lowest miss point estimate under hectic-day saturation (\RHMiss\%), though its baseline miss difference from CADR is not significant in the 30-seed test; its lead becomes significant only at the scarcer $\times 0.75$ availability scale. Its Phase~3 wait has no significant paired difference from SPT's and is in the lowest tier (\RHWait~min, $p=\RHSPTwaitP$; at monthly scale its wait no longer leads, sitting within about two minutes of SPT's, CADR's and SPT+Rescue's). Its load-gated anticipation remains a genuine gain below saturation, cutting surge-day miss from \ResSurgeOffMiss\% to \ResSurgeOnMiss\% and placing RH, with EDF and the Adaptive hybrid, below CADR on surge-day miss (day-type sweep, no paired test), at a higher surge-day wait (\SensSrgRHwait~min). Its high tardiness (\RHtard~min Phase~3), exceeding even the FIFO baseline, means that the deadlines it does miss are missed by a large margin, so RH suits a service-level agreement that counts breaches regardless of lateness, in a regime where per-type availability is known to be scarce in the assumed order.
  \item \textbf{Use SPT as a strong low-wait baseline under heavy load ($\rho>1$).} SPT achieves a lowest-tier wait, with no significant paired wait difference from RH (\SPTwait~min, \SPTFIFOwaitDiffPct\% below FIFO's \FIFOwait~min), though EDF now achieves significantly lower miss. At monthly scale: \PhFourSPTwait~min wait, \PhFourSPTmiss\% miss.
  \item \textbf{Deploy SPT+Rescue as the conservative fallback when tight-deadline compliance is a strict requirement.} SPT+Rescue's laxity-rescue mechanism limits starvation even under saturation (\RescueMiss\% miss, notably lower than SPT's \SPTmiss\%) at marginally higher wait (\RescueWait~min vs SPT's \SPTwait~min).
  \item \textbf{Use EDF at lower loads ($\rho<1$) with mixed deadline distributions.} Among the baseline policies in the deadline sweep, EDF achieves significantly lower miss than SPT at low tight-deadline fractions ($\leq 30\%$; \DlTenEDFmiss\% vs SPT's \DlTenSPTmiss\% at 10\% tight, $\DlTenSPTEDFmissP$), and \DlTenFIFOmiss\% for FIFO at the same point. Switch to SPT only at high tight fractions ($\geq 70\%$, \DlSeventySPTmiss\% vs EDF's \DlSeventyEDFmiss\% at 70\%, $\DlSeventySPTEDFmissP$); at the intermediate $f_{\text{tight}}=50\%$ point the paired test does not resolve their difference ($\DlFiftySPTEDFmissP$), so neither ordering is preferable there. Under hectic saturation, SPT dominates EDF on wait; EDF achieves significantly lower miss (\EDFmiss\% vs SPT's \SPTmiss\%, $p<0.001$).
  \item \textbf{To reduce the residual monthly miss rate further, increase $C_{\max}$.} The miss rate at $C_{\max}=5$ is the price of online, non-preemptive, stochastic operation under the concurrency cap, and no ordering policy removes it; extra slots absorb provisioning holding time and variance (Section~\ref{ssec:capacity-lower-bound} shows that a miss-free schedule exists on the expected-service instances, so the gap is not forced by capacity alone). The $C_{\max}$ sensitivity analysis (Section~\ref{ssec:concurrency-limit-sensitivity}) confirms that, at hectic-day load, $C_{\max}=5$ sits at the inflection point of the saturation curve: the gain from increasing to 7 is substantially smaller than the gain from 3 to 5, at $C_{\max}=10$ every scheduler, FIFO included, reaches zero miss on the hectic day. Higher concurrency (available through RunPod's balance-tiered quotas or on request) is the primary lever for reducing both waiting time and miss rate under saturation. The rescue threshold sensitivity (Section~\ref{ssec:rescue-threshold-sensitivity}) additionally shows that SPT+Rescue miss varies by \ThrMissSpread\,percentage points across the swept range $\theta \in \{60, 180, 300, 600, 1200\}$\,s ($n=10$ seeds); $\theta=60$\,s is the weakest setting and the lowest miss, at $\theta=180$\,s, lies within the seed noise and was not retested at $n=30$, so the default $\theta=600$\,s is acceptable and values from about 180\,s upward sit in a narrow band. The deadline-aware advice for $\rho<1$ above (EDF relative to SPT) concerns the baseline policies only; CADR remains the default for saturated and monthly operation, while on surge days EDF, the Adaptive hybrid and RH have lower miss point estimates (day-type sweep, no paired test).
\end{enumerate}

\subsection{Choosing Between RH and CADR}
\label{ssec:choosing-between-rh-and-cadr}

RH and CADR are the two strongest schedulers in this study, and an operator deploying one of them faces a single, well-defined choice. Both have low deadline-miss point estimates in Phase~3 (RH \RHMiss\%, CADR \CADRmiss\%); their paired difference is not significant ($p=\RHCADRmissP$), but the 95\% interval for RH minus CADR [\RHCADRmissMDLo, \RHCADRmissMDHi] percentage points does not establish equivalence. They remain close on the Phase~4 point estimates (RH \PhFourRHMiss\%, CADR \PhFourCADRMiss\%, no paired test at that phase), and at Phase~3 hectic saturation both post lower miss point estimates than SPT+Rescue (a moderate effect, no pre-specified paired test), though neither dominates it on wait: RH's wait advantage over SPT+Rescue does not survive statistical correction, and CADR's wait is essentially tied with it. Crucially, in Phase~3 both lie on the (miss rate, tardiness) Pareto frontier, at its low-miss end (its low-tardiness end is held by the Adaptive hybrid, \AdaptiveTard~min at \AdaptiveMiss\% miss), and neither dominates the other there (at Phase~4 RH has the lower miss point estimate and CADR the lower tardiness point estimate, with no paired test), because they diverge on one quality axis, the \emph{severity} of the deadlines they do miss. RH promotes a tight job only in its last service-time window, so a tight job that misses it is demoted behind the loose backlog and missed by a large margin (mean tardiness \RHtard~min in Phase~3); CADR's earlier critical-ratio promotion cuts that severity by \CADRtardReduction\% (\CADRtard~min) while the baseline miss difference is not resolved by the paired test. The miss comparison is itself regime-dependent: RH's lead over CADR is significant only at the scarcer $\times 0.75$ availability scale, and CADR's miss is lower at $\times 1.25$ and under the data-ordered per-type ranking, the ranking the snapshot offer shares support (Section~\ref{ssec:availability-model-sensitivity}). CADR is therefore the recommended default: it has an inconclusive baseline miss comparison with RH, has close monthly miss point estimates with lower tardiness, and leads under eased availability and the data-ordered ranking, and cuts lateness substantially. RH remains the right choice in a narrower setting, the scarce baseline availability regime under a service-level agreement (SLA) that counts breaches regardless of lateness, as summarised in Table~\ref{tab:rh_vs_cadr}.

\begin{itemize}
  \item \textbf{Default to CADR, in particular when the penalty scales with lateness.} If late jobs incur escalating cost, for example hourly late fees or a downstream pipeline stage (compositing, review) that degrades the longer a frame is late, CADR is clearly preferable: it holds the miss count within one percentage point of RH (\CADRmiss\% vs \RHMiss\% at Phase~3, slightly above it at Phase~4) while cutting mean lateness by \CADRtardReduction\% (\CADRtard~min vs \RHtard~min Phase~3; \PhFourCADRTard~min vs \PhFourRHtard~min Phase~4). On the tested tardiness-weighted grid, CADR is the lower-composite policy despite missing marginally more deadlines at Phase~3. It is also the default when lateness does not matter, because the baseline paired miss test does not resolve a difference from RH and it is the more robust of the two: under the data-ordered ranking it leads on miss and RH is significantly worse.
  \item \textbf{Specialise to RH when the penalty is per-breach, lateness-independent, and the scarce baseline availability regime holds.} If the SLA counts each missed deadline equally regardless of how late the job finishes (for example, a fixed credit per breached rush order) and per-type availability is known to be scarce in the assumed order, RH is a defensible choice, though not a demonstrably better one on this criterion: it has the lowest hectic-saturation miss point estimate (\RHMiss\% Phase~3), but this is not a significant difference from CADR's in the baseline paired test ($p=\RHCADRmissP$), its lead becomes significant only in the scarcer $\times 0.75$ availability regime, and it reverses under the data-ordered availability ranking (Section~\ref{ssec:availability-model-sensitivity}); it also has a monthly miss near-identical to CADR's (\PhFourRHMiss\% vs \PhFourCADRMiss\% Phase~4, no paired test) together with best-tier queue wait at hectic saturation (\RHWait~min, tied with SPT, $p=\RHSPTwaitP$; at monthly scale RH's wait trails SPT's by under two minutes and is slightly below CADR's and SPT+Rescue's). The large tardiness of its few misses carries no extra penalty under such an SLA.
\end{itemize}

\begin{table}[htbp]
\centering
\caption[RH vs CADR decision guidance]{Decision guidance for the two proposed schedulers: CADR is the default, and RH is the specialist for the scarce baseline availability regime. The choice is a service-quality choice (robustness and miss severity versus the lowest breach point estimate in one regime); it is not a cost trade-off, because operational cost differs by only about \CostSpreadPct\% across all ten policies under saturation (RH \$\RHcost{} vs CADR \$\CADRcost{} per 24-hour run, under 1\% apart).}
\label{tab:rh_vs_cadr}
\begin{singlespace}
\small
\begin{tabular}{>{\raggedright\arraybackslash}p{4.3cm}l>{\raggedright\arraybackslash}p{6.5cm}}
\toprule
\textbf{Operator priority} & \textbf{Choose} & \textbf{Decisive metric} \\
\midrule
Default (no strong prior on the SLA or on the availability regime) & CADR & Phase~3 miss difference from RH not significant ($p=\RHCADRmissP$), equivalence untested; close at Phase~4 (\PhFourCADRMiss\% vs \PhFourRHMiss\%); lower at $\times 1.25$ and under the data-ordered ranking \\ \addlinespace[4pt]
Lateness penalised (hourly late fees, latency-sensitive downstream stages) & CADR & \CADRtardReduction\% lower tardiness (\CADRtard{} vs \RHtard~min) \\ \addlinespace[4pt]
Per-type availability ranking uncertain & CADR & Leads on miss under the data-ordered ranking (\RevCADRmiss\% vs RH \RevRHmiss\%, $\RevRHvCADRp$) \\ \addlinespace[4pt]
Fewest breached deadlines in the scarce baseline regime (per-breach SLA, lateness-independent) & RH & Lowest Phase~3 miss point estimate (\RHMiss\%, with no significant paired difference from CADR); lead significant only at $\times 0.75$ \\ \addlinespace[4pt]
Minimise queue wait at hectic saturation & SPT & Best-tier Phase~3 wait at $C_{\max}=5$ only (\SPTwait~min, with no significant paired wait difference from RH's \RHWait~min; at $C_{\max}=10$ RH has the highest wait, \CmaxTenRHwait~min against \CmaxTenOtherWaitLo{} to \CmaxTenOtherWaitHi~min for the other policies); at monthly scale RH, CADR and SPT+Rescue sit within about two minutes of one another (\PhFourRHWait{}, \PhFourCADRwait{} and \PhFourRescueWait~min respectively, no paired policy test) behind SPT \\ \addlinespace[4pt]
Lowest operational cost & neither & ${\sim}\CostSpreadPct\%$ spread; quality is governed by $C_{\max}$ (its cost effect was not measured), not by the scheduler \\
\bottomrule
\end{tabular}
\end{singlespace}
\end{table}

\textbf{Service quality versus cost.} As Table~\ref{tab:rh_vs_cadr} summarises, the small hectic-day cost difference is not the basis for choosing between RH and CADR. Under the saturated hectic-day regime ($\rho>1$), every work-conserving policy keeps the GPU fleet at the same high utilisation, so total operational cost is nearly fixed across all schedulers: the spread between the cheapest policy (SPT, \$\SPTcost{} per 24-hour run) and the most expensive (LCF, \$\LCFcost{}) is only about \CostSpreadPct\%, a marginal, near-noise gap, and RH (\$\RHcost{}) and SPT+Rescue (\$\RescueCost{}) sit inside that band. Between the two proposed schedulers the difference is under 1\% (RH \$\RHcost{} vs CADR \$\CADRcost{} per 24-hour run, under one cent), negligible beside the miss-count and tardiness differences. On the hectic-day point estimates, cost differences between RH and CADR are small relative to the service-quality differences; Figure~\ref{fig:cost-comparison} in Chapter~\ref{chap:experiments} visualises this compressed cost range for both the hectic-day and the monthly horizon. The genuine service-quality-versus-cost trade-off in this system is not located in the scheduler at all but in the concurrency limit $C_{\max}$: raising $C_{\max}$ provisions more parallel capacity, lowering both wait and miss; its effect on spend was not measured here, and because only realised execution is billed it can change spend only through the GPU-type mix, not through idle capacity (Section~\ref{ssec:concurrency-limit-sensitivity}). An operator therefore compares service quality at similar, not identical, hectic-day spend when choosing the scheduler (CADR by default and for lower lateness, RH for the lowest breach point estimate in the scarce baseline regime); the monthly cost spread is wider, from RH's mean of \$48.56 to FIFO's \$54.36 across the ten replications, so \emph{fixed cost} does not describe the monthly horizon. Choosing $C_{\max}$ changes the provider quota tier or cap, with a spend effect not measured in the concurrency sweep.

\section{Limitations and Future Research}
\label{sec:limitations-and-future-research}

\subsection{Limitations}
\label{ssec:limitations}

The following limitations bound the scope of the conclusions.

\textbf{Simulation-only evaluation.} Every result reported here is produced by the discrete-event simulator; no scheduler was run against the live RunPod platform, so no figure in this thesis is a field measurement. The point bites hardest on the availability model, whose stock-status probabilities and provisioning delays are assumed rather than measured (see the availability-model limitation below), so the absolute miss, wait and tardiness levels quoted throughout are outputs of that model and not observed platform behaviour. What the study claims is the \emph{relative ordering} of the ten policies under one transparent, fully specified model, and the deployment rules of Section~\ref{sec:theoretical-implications} should be read on those terms: an operator adopting them should expect the separation between deadline-aware and deadline-blind policies to carry over, but neither the numbers nor the order within the leading group, which passes from RH to CADR, SPT+Rescue, EDF and the Adaptive hybrid (all at zero miss) at $\times 1.25$ availability, and to CADR and SPT+Rescue when its per-type ranking follows the snapshot offer shares (Section~\ref{ssec:availability-model-sensitivity}). Live validation on the platform is the fourth item of the future work below, and it is the item that would turn these rules from model-supported into field-supported.

\textbf{Arrival-rate knowledge.} RH's anticipation gate compares an assumed arrival rate with a fixed switching point. On the hectic day an assumed rate below about \PriorGateLambda~jobs/s (more than about \PriorGateShortfall\% under the true rate) switches the gate the wrong way and more than doubles RH's miss rate (Section~\ref{ssec:anticipative-rolling-horizon-scheduler}); a conservative high estimate fails safe.

\textbf{Single-source workload model.} The job model rests on the submission history of one studio's rendering pipeline, and two of its ingredients inherit that pipeline's shape directly. The complexity levels $\kappa$ are assigned by render-time terciles of that history, which makes $\kappa$ an ex-post label that the simulator treats as known at arrival (an oracle tercile; a bracketing check with noisy and uninformative labels in Section~\ref{ssec:label-noise-and-gpu-speed-assumptions} leaves RH and CADR with the lowest miss point estimates but at higher miss rates than under the oracle label, and no real classifier has been measured), and the boundary between low, medium and high complexity is that studio's asset mix rather than a general one; and its 14:00 to 17:00 submission peak falls in the diurnal model's lowest-availability band, while the hectic and surge daily job counts are stress loads above anything in that history (the simulated hectic day itself opens at 08:00 and its arrival window lies in the morning band, Section~\ref{ssec:platform-and-data}). A studio with heavier scenes, a different render farm, or an overnight submission habit would shift both, changing the tight/loose mix a scheduler faces and when the queue builds. The deadline split (20\% tight / 80\% loose) is a modelling assumption on top of that, swept in Section~\ref{ssec:deadline-tightness-sensitivity} but never observed. Reproducing the comparison on a second studio's history is the cheapest available test of how far the ordering generalises.

\textbf{Availability model simplification.} The availability model is not calibrated to measured provisioning behaviour. The author's pricing-API record (Section~\ref{ssec:platform-and-data}) reports the fleet types' Secure Cloud stock as almost never High, so the baseline high-stock probabilities $\beta_g$ are assumed fulfilment probabilities rather than observed frequencies, and the provisioning-delay ranges were not measured; the record supports the diurnal shape of availability only qualitatively, and its per-type offer shares rank the four fleet types in the reverse of the assumed $\beta_g$ order. The availability-model sensitivity sweep (Section~\ref{ssec:availability-model-sensitivity}) bounds this concern: across the scarcity side of the operable range ($\times 0.75$ to $\times 1.00$ of the baseline high-stock probabilities) the scheduler comparison is stable, so on that side the conclusions do not hinge on the exact calibration point. The per-type ranking matters more: under the data-ordered permutation that matches the offer shares, CADR (\RevCADRmiss\%) and SPT+Rescue (\RevRescueMiss\%) lead and RH (\RevRHmiss\%) is significantly worse than CADR, so RH's miss lead is conditional on the assumed ranking. At $\times 0.50$ and below the platform becomes inoperable for every policy alike, and at $\times 1.25$ the eased regime lets several policies reach zero miss, so the RH advantage is a property of the baseline per-type ranking and scarcity-and-saturation regime rather than a universal ordering.

\textbf{Data-versus-model gaps.} Four modelling choices depart from the production system the data came from (Section~\ref{ssec:platform-and-data}). (1) Production ran on a RunPod Serverless endpoint, but the simulator prices each slot at the Pods \texttt{securePrice}, lower than the Serverless Flex rates by a factor of about 1.5 to 2.5; under Serverless pricing the RTX~3090 and RTX~A5000 cost the same, so absolute costs and the RTX~3090 cost premium are specific to Pods pricing. (2) The service time is render time only: the scene-import phase (mean 10.8\,s, about one seventh of the mean render time) is excluded from occupancy and cost, so both are optimistic by roughly that proportion. (3) The hectic (950) and surge (730) job counts are stress-load assumptions: the busiest recorded day holds 415 jobs, so the saturated-regime conclusions describe a heavier workload on this fleet than the recorded one. (4) $C_{\max}=5$ is RunPod's default Serverless worker quota, adopted as a baseline; the 393 logged jobs that ran while more than five jobs were in flight are excluded from the calibration data.

\textbf{Calibration sample and fleet composition.} The calibration set includes 203 jobs tagged as development runs and one untagged job (3.3\% of the 6,234). Excluding them leaves the GPU speed ordering unchanged, but the per-stratum fit was not rerun on production rows alone. The experimental fleet is also fixed at two RTX~3090 slots and one each of the RTX~A5000, A4500 and A4000, matching the four types present in the data. A six-mix sensitivity check at ten seeds (Section~\ref{ssec:fleet-composition-sensitivity}) shows that the lead is composition-dependent: RH and CADR are the two lowest-miss policies in four of the six mixes and every deadline-aware policy stays below FIFO in all six, but the A5000-heavy mix places CADR and SPT+Rescue first and second with RH third by a difference the paired test does not resolve ($p=0.282$, pooled $d=0.287$), and the A4000-heavy mix places CADR first at 0.57\% miss against RH's 5.96\% ($p<0.001$, pooled $d=6.21$). The sweep is a sensitivity check on one workload, not a second studio trace, and cross-mix service draws differ in dispatch order even though the anchor mix reproduces the first ten Phase~3 seeds exactly. A second studio trace remains the open generalisation test.

\textbf{Arrival process and service-time fit.} Arrivals are Poisson within each day type, so inter-arrival times are memoryless and the simulated stream does not reproduce the measured hourly arrival profile; a non-homogeneous Poisson process or a trace-driven replay is left to future work. The per-stratum service-time model is moment-matched (mean and log-space dispersion) and no formal goodness-of-fit statistic is reported. The distributional robustness check is the mean-based rerun in Section~\ref{ssec:platform-and-data}, which leaves the RH and CADR ordering unchanged, but the fit itself is not tested.

\textbf{Baseline coverage.} The comparison now includes HEFT (which, for independent jobs, reduces to longest-average-job-first ordering), min-min, max-min, and a shallow tabular bandit with a random-arm control (Section~\ref{ssec:classical-and-learned-baselines}). HEFT and max-min coincide on all 30 seeds and all six metrics and are reported as one policy; the paired test does not separate MinMin's wait from RH's, but MinMin misses significantly more deadlines; and the shallow bandit does not beat FIFO, while its random-arm control records a lower miss rate than the trained bandit. These results close the classical-baseline gap for the dispatching problem studied here. A metaheuristic over the static offline relaxation, and a stronger learned policy, remain open; deep reinforcement learning in particular is not evaluated (see the reinforcement-learning limitation below).

\textbf{Non-preemptive assumption.} CGRSP models non-preemptive scheduling. Some cloud platforms do permit early termination (spot preemption). Extending the formulation to preemptive scheduling would require tracking partial completion and migration overhead.

\textbf{Single-job-per-frame model.} This study models each rendering frame as a single indivisible job: one job is one complete frame, not a pixel, a sample, a tile or a sub-frame, and several frames are not combined into one job. In practice, some renderers support tile-based rendering (splitting a frame into spatial tiles rendered in parallel on multiple GPUs), which would require a DAG-structured job model.

\textbf{Sparse strata.} The 6,234-job calibration dataset is real but unevenly distributed across the twelve $(\kappa, g)$ strata: the RTX~A5000 dominates (4,680 jobs) and the RTX~A4000/A4500 are sampled almost entirely at high complexity. Three of the twelve strata are too sparse for a direct per-stratum fit and instead use a relative-performance estimate anchored on the same GPU's well-sampled cells: low-complexity A4000 (no jobs), low-complexity A4500 (no jobs), and medium-complexity A4000 (two jobs). These cells receive \SparseExposureLo\% to \SparseExposureHi\% of synthetic dispatches across the six focal policies at the ten-seed baseline, so their effect cannot be dismissed from the historical sample counts. The direct mean and dispersion perturbations in Section~\ref{ssec:uncertainty-in-estimated-execution-times} change absolute miss rates and can reverse the point-estimate ordering of RH and CADR. They are stress brackets rather than uncertainty intervals inferred from data; collecting matched renders for these cells remains necessary to reduce that uncertainty. The A4000 means are also not scene-matched: on the only two scenes rendered on both the A4000 and the A5000 the ratio is 1.165, and the scene-matched arm raises every policy's miss rate while RH and CADR stay lowest (Table~\ref{tab:kappa_speed_sensitivity}, Section~\ref{ssec:label-noise-and-gpu-speed-assumptions}).

\textbf{No nontrivial optimality bound.} The CGRSP instance sizes studied ($|\calJ| \approx 950$, $C_{\max}=5$, $n=30$ replications) make exact integer linear programming intractable at the required evaluation scale; no polynomial-time optimal algorithm is known for the general heterogeneous parallel machine scheduling problem with deadlines~\cite{pinedo2016scheduling}. The capacity analysis of Section~\ref{ssec:capacity-lower-bound} certifies only that the offline-relaxed optimum on the expected-service instances is zero misses, so it provides only a trivial (zero) benchmark, not a meaningful gap against which the policies could be scored. Scheduler performance is therefore evaluated relative to the FIFO baseline and relative to each other. SPT's near-M/G/1-optimal mean-wait behaviour under saturation provides an indirect argument for near-optimality in the wait dimension; no such bound is available for miss rate.

\textbf{Deep reinforcement learning not evaluated.} Deep-RL policies (DQN, policy gradient) were not included in the comparison for three reasons: (1) training overhead: deep-RL schedulers must be trained on simulated workloads before deployment~\cite{chen2023gpu}, a cost the heuristics do not incur; (2) the scheduling environment changes frequently as RunPod updates GPU types and pricing, requiring periodic retraining; and (3) interpretable deployment rules are more operationally actionable for rendering studios than black-box policies. Deep RL remains a direction for future work (see the reinforcement-learning item below). The shallow tabular bandit in Section~\ref{ssec:classical-and-learned-baselines} is a separate baseline, not a deep-RL policy; its null result is reported there.

\noindent\textbf{Reproducibility.} The simulator, job data, generated tables and figure sources accompany this thesis. The pipeline is deterministic under a fixed hash seed, and a validation script checks every reported number against the generated macros and the recorded results. The repository is available to the examination committee and will be deposited publicly on acceptance.

\Needspace{6\baselineskip}
\subsection{Future Research}
\label{ssec:future-research}

Thirteen directions for future work follow from the limitations above and from the scope of the present study:

\begin{enumerate}
  \item \textbf{Preemptive scheduling extension.} Incorporating checkpoint-and-resume (available in some cloud renderers) would allow modelling preemption events and evaluating EDF's theoretical guarantees in a preemptive setting.

  \item \textbf{Multi-objective Pareto exploration.} The current Pareto analyses are limited to the (miss, wait) and (miss, tardiness) planes. A full three-objective Pareto analysis (waiting time, miss rate, cost) using NSGA-II or similar multi-objective evolutionary algorithms would provide a richer trade-off characterisation.

  \item \textbf{Tile-based rendering job model.} Extending CGRSP to DAG-structured jobs (where a frame is decomposed into tiles) would enable scheduling research on job-level parallelism, relevant to real-time rendering pipelines, and would connect the fleet dispatcher to kernel-level scheduling frameworks~\cite{chien2018graph,chang2020dependent}.

  \item \textbf{Real-platform integration.} Deploying the scheduler as a middleware layer on top of the RunPod API (using the platform's job submission endpoint) would enable live validation of the simulation results in production conditions. In particular, measuring per-type fulfilment probabilities would settle whether the baseline or the data-ordered availability ranking holds, which decides whether RH's specialist regime occurs in practice or CADR's lead under the data-ordered ranking applies.

  \item \textbf{Second-studio replication.} Reproducing the comparison on a second studio's submission history would test how far the ordering generalises beyond the one trace studied here; the six-mix fleet-composition sweep is reported in Section~\ref{ssec:fleet-composition-sensitivity}, and a second workload remains the open test.

  \item \textbf{Non-homogeneous arrivals and trace-driven replay.} Replacing the within-day Poisson stream with the measured hourly arrival profile, or replaying the recorded arrival trace directly, would remove the strongest remaining arrival-process assumption.

  \item \textbf{A stronger learned policy.} The classical baselines are now in the comparison (Section~\ref{ssec:classical-and-learned-baselines}); what remains is a learned policy stronger than the shallow bandit tested there, which is the subject of item~12 below.

  \item \textbf{RH plus CADR hybrid.} RH's timeline determines which jobs are doomed, while CADR promotes at-risk jobs earlier with a critical ratio. A hybrid that keeps RH's timeline and widens its promotion window could test whether the miss-count versus lateness trade-off can be escaped rather than only spanned.

  \item \textbf{Newer GPU types.} The experimental fleet is limited to four GPU types available on RunPod Secure Cloud in March 2026. Extending to more recent GPU generations (H100, RTX~5090) and investigating the effect of higher $C_{\max}$ limits (available through RunPod's balance-tiered quotas or on request) beyond the swept range ($C_{\max}$ up to 10), such as the higher balance-tier quotas RunPod offers up to 60 workers, would extend the concurrency analysis and allow measuring its cost effect.

  \item \textbf{Deadline distribution assumptions.} The deadline distribution used in this study (20\%/80\% tight/loose split) is a modelling assumption. A study using real production schedules from a rendering studio would validate the choice of threshold and quantify the benefit of EDF in practice.

  \item \textbf{Load-adaptive threshold policy.} The threshold sensitivity study in Section~\ref{ssec:rescue-threshold-sensitivity} shows SPT+Rescue performance is relatively insensitive to $\theta$ across the swept range (miss within \ThrMissSpread\,percentage points; the sweep minimum is at $\theta=180$\,s, but the within-band differences are comparable to the sampling noise at $n=10$). Future work could explore a load-adaptive threshold policy that adjusts $\theta$ dynamically based on current queue pressure and GPU utilisation.

  \item \textbf{Reinforcement learning scheduling policies.} The shallow tabular bandit of Section~\ref{ssec:classical-and-learned-baselines} is a deliberately limited learned baseline: it selects among four fixed arms with episode-level credit assignment and does not beat FIFO, so it neither supports nor rules out deep reinforcement learning. Deep RL methods (policy gradient and DQN variants) have shown promise for cluster and GPU job scheduling~\cite{mao2016resource,chen2023gpu}. Applying RL to CGRSP (with a state space encoding queue laxity distribution, GPU stock status, and time-of-day) could discover non-myopic dispatch policies that outperform the greedy heuristics studied here. The CGRSP simulator provides a ready training environment.

  \item \textbf{Automated heuristic discovery.} The CGRSP policies are hand-designed. LLM-driven heuristic search~\cite{ye2024reevo,cetinkaya2025discovering}, which has already been used to extend due-date dispatching rules, could use the CGRSP simulator as its evaluation harness to search for dispatch rules that improve on RH and CADR while remaining interpretable.
\end{enumerate}

\section{Concluding Remarks}
\label{sec:concluding-remarks}

Scheduling 3D rendering over on-demand GPU fleets with marketplace-driven availability has not previously been formulated or evaluated: rendering studios using RunPod or similar platforms have had no principled guidance on which scheduling policy to deploy. This thesis provides that formulation and the accompanying empirical evidence.

Within the fixed concurrency budget ($C_{\max}=5$), the heterogeneous tight/loose deadline formulation, and the hectic overload regime ($\rho>1$), RH has the lowest deadline miss point estimate (\RHMiss\%, significantly below SPT+Rescue's, $p<0.001$, though not significantly different from CADR's, $p=\RHCADRmissP$) and best-tier wait (\RHWait~min, tied with SPT), at the cost of high mean tardiness (\RHtard~min, above FIFO's). CADR achieves near-RH miss (\CADRmiss\%) at substantially lower tardiness (\CADRtard~min). The six-mix fleet sweep (Section~\ref{ssec:fleet-composition-sensitivity}) shows that RH's lead is composition-specific: RH is the lowest-miss point estimate in four of the six mixes, while CADR is never more than about one percentage point above the mix minimum and leads in the other two, so CADR is the more robust default across fleet compositions. SPT achieves a lowest-tier wait (\SPTwait~min, tied with RH) but a significantly higher miss rate than EDF (\SPTmiss\% vs \EDFmiss\%). SPT+Rescue is the conservative fallback: a lower miss rate and lower wait than FIFO, with its wait premium over SPT. EDF beats SPT on miss but has a higher mean wait than FIFO (not significant after correction). FIFO/LCF/Balanced incur \FIFOmiss\% to \Balancedmiss\% miss and should not be deployed when deadline compliance matters. CADR is the recommended default: its baseline paired miss difference from RH is not significant (equivalence untested), it has close monthly miss point estimates with lower tardiness, and it leads under eased availability and the data-ordered per-type ranking. RH remains a genuine contribution and the specialist choice for the scarce baseline availability regime when only the count of missed deadlines matters, with load-gated anticipation that cuts surge-day miss below saturation. Increasing $C_{\max}$ beyond the default budget of five is the primary mechanism to reduce residual miss rates at saturation.

} % end \sloppy
